\chapter{The File-System Tree}
\label{ssec:fs-tree-internal}

The chapter contains the single most consequential reverse-engineering
finding in the project so far. The finding is small; its absence cost
a fixture and almost cost the v1 contract; its presence makes the
parser correct on every FS tree dense enough to grow past one leaf
node.

\section{The Finding, Stated Once}

\textbf{Observation.} An FS-tree internal-node entry has an 8-byte
value. That value is a \emph{virtual object identifier}, not a
physical block address. To read the child node, the walker must
resolve the value through the \emph{volume} OMAP at the scan XID.
A walker that reads the value as a paddr will, on any FS tree of
depth $\geq 1$, read random blocks.

This is the rule. The rest of the chapter develops the consequences
and the empirical evidence.

\section{Why This Is The Easy Mistake}

The Apple reference documents that FS-tree objects live in the volume's
virtual address space. It does not foreground that
\emph{internal-node child references are members of the same address
space}. The reference reads naturally as "the internal value is a
pointer into the same area as physical extents," and many B-tree
formats indeed encode internal child references as direct physical
offsets.

The converged open-source code is unanimous in the right direction,
but the unanimity is implicit. Each reader hands its B-tree walker an
"object resolver" that internally does the OMAP lookup, and the
walker code does not need to think about it. A reader who reimplements
the walker without an equivalent resolver --- which is what the
project's first Rust walker did --- will not notice the lapse on a
small fixture, because a small enough FS tree fits in a single leaf
node and has no internal entries to follow.

The project missed this through the EX-10, EX-11, and EX-12 fixtures.
Each was small enough that the FS tree had no internal level, and the
walker was never asked to dereference an internal value. EX-14
finally produced a fixture dense enough to need an internal level,
and the parser broke.

\section{The Probe That Found It}

\textbf{Observation.} EX-14 reported
\texttt{APFS object validation failed: checksum mismatch at block
1031} and aborted before publishing a selected checkpoint. The probe
\texttt{fsck\_apfs -n} returned clean; \texttt{go-apfs identitydump}
succeeded on the same container. The image was not malformed.

EX-15 replayed the upstream context in Python, validating every
candidate against the strict rules of chapters~3 and 4. Every NXSB
candidate validated. Every checkpoint map block validated. The
container OMAP, volume superblock, and volume OMAP all validated.
The FS-tree root validated. The walker then dereferenced an internal
entry whose 8-byte value was \texttt{1031}, read block 1031, and
encountered 4096 bytes of zero --- an unallocated block --- whose
Fletcher-64 mismatched.

The volume OMAP, queried for \texttt{oid=1031} at the scan XID, returned
\texttt{paddr=505}. Block 505 was a valid FS-tree internal node. The
walker had been asked to follow an OID, and had instead read the OID's
numeric value as a paddr.

\begin{figure}[h]
\centering
\begin{tikzpicture}[
  every node/.style={font=\small},
  node distance=0.55cm and 1.2cm,
  rootnode/.style={draw, fill=blue!10, minimum width=4.2cm, minimum height=1.15cm,
                   align=center, font=\footnotesize, rounded corners=2pt},
  omapnode/.style={draw, fill=orange!18, minimum width=2.6cm, minimum height=1.0cm,
                   align=center, font=\footnotesize, rounded corners=2pt},
  child/.style={draw, fill=green!18, minimum width=2.0cm, minimum height=0.85cm,
                align=center, font=\footnotesize},
  bad/.style={draw, fill=red!12, minimum width=2.0cm, minimum height=0.85cm,
              align=center, font=\footnotesize, dashed},
  rule/.style={font=\scriptsize, align=center},
  good/.style={-Stealth, thick, green!50!black},
  wrong/.style={-Stealth, thick, red!70!black, dashed},
]
  \node[rootnode] (root)
    {FS-tree root \quad (paddr 501)\\
     \scriptsize \texttt{level=1, nkeys=2}, internal values \texttt{1030},
     \texttt{1031}};

  \node[bad, below left=1.5cm and -0.3cm of root] (b1031)
    {block 1031\\\scriptsize unallocated, zeros};

  \node[omapnode, below right=1.5cm and -0.3cm of root] (omap)
    {volume OMAP\\\scriptsize lookup\\\scriptsize \texttt{(oid, scan\_xid)}};

  \node[child, below=0.9cm of omap] (block505)
    {block 505\\\scriptsize valid child};

  \draw[wrong] (root.south west) -- (b1031.north)
    node[midway, left, rule]
      {wrong:\\read value\\as paddr};

  \draw[good] (root.south east) -- (omap.north)
    node[midway, right=2pt, rule]
      {right:\\value is a\\virtual OID};

  \draw[good] (omap.south) -- (block505.north)
    node[midway, right=2pt, rule]
      {\texttt{1031 $\to$ 505}};
\end{tikzpicture}
\caption{The headline finding. An FS-tree internal-node value is a
  virtual OID; the walker must consult the volume OMAP at the scan
  XID before reading the child. Reading the value as a paddr lands on
  an unallocated block.}
\end{figure}

The fix is structural. \texttt{walk\_fs\_node} now takes the volume
OMAP resolver and performs an \texttt{(oid, scan\_xid)} lookup for
every internal entry before reading the child. A missing OMAP mapping
returns a typed \texttt{InvalidObject} hard stop naming the OID.

Two synthetic regression tests in the Rust crate close the bug:

\begin{enumerate}
  \item \texttt{fs\_tree\_internal\_value\_is\_virtual\_oid\_resolved\_via\_omap}
    builds an 8-block synthetic image where the FS-tree root's
    internal entry has value \texttt{1030}. With the fix, the walker
    follows OMAP to the correct paddr and reads the leaf record.
    Without the fix, the walker short-reads on block 1030.
  \item \texttt{fs\_tree\_internal\_oid\_missing\_from\_omap\_is\_hard\_stop}
    asserts that an internal value pointing at an OID with no OMAP
    mapping returns a typed error mentioning the missing OID rather
    than silently skipping or attempting a paddr fallback.
\end{enumerate}

These tests now run on every commit; a future regression that
reintroduces the bug would be caught at the synthetic level rather
than at the next fixture variant.

\section{The Object-Header Validation Detail}

Two further details follow from the FS-tree root being virtual.

\textbf{Observation.} The FS-tree root is reached through the volume
OMAP, so on disk it carries a virtual \texttt{obj\_phys\_t} header.
The storage flags are \texttt{0x0000\_0000} (virtual), not
\texttt{OBJ\_PHYSICAL}.

\textbf{Observation.} A virtual object's \texttt{o\_oid} field does not
equal the block's physical address; the OID is the virtual identity,
the paddr is independent. The validator therefore must \emph{not}
require \texttt{o\_oid == paddr} for virtual objects. The right check
is \texttt{o\_oid == omap\_lookup.oid}: the OID the walker resolved
through the OMAP must be the OID the on-disk header advertises.

The corollary is that the project's object-header validator has two
modes: \texttt{ExpectedStorage::Physical}, which requires
\texttt{o\_oid == paddr}, and \texttt{ExpectedStorage::Virtual}, which
does not. Every reader picks one at the call site. The fail-closed
rule is preserved: a block validated under the wrong mode is
rejected.

\section{Required Record Families}

Narrow v1 needs five families of FS-tree records and one xfield
family.

\begin{description}
  \item[\texttt{DIR\_REC}.] Directory entries. Key carries the parent
    inode's file ID and the name hash; value carries the child file
    ID and entry type. The directory structure is built by walking
    DIR\_REC records.
  \item[\texttt{INODE}.] File, directory, symlink, and other entry
    metadata. Carries the parent ID, link count, mode, internal
    flags, and an xfield tail.
  \item[\texttt{XATTR}.] Extended attributes. v1 cares about three
    specifically: \texttt{com.apple.fs.symlink} (the symlink target),
    \texttt{com.apple.decmpfs} (the compression metadata header), and
    \texttt{com.apple.ResourceFork} (out of v1 scope but tracked).
  \item[\texttt{SIBLING\_LINK}, \texttt{SIBLING\_MAP}.] Hard-link
    unification. A directory entry that participates in a hard-link
    group carries a sibling ID xfield; SIBLING\_MAP closes the
    mapping from sibling ID to file ID.
  \item[dstream fields.] Logical size sources. Carried either inside
    INODE's xfield tail as \texttt{INO\_EXT\_TYPE\_DSTREAM} or as a
    dedicated \texttt{j\_dstream\_id\_val\_t} record.
\end{description}

Other families are decoded but not used for v1 product rows:
\texttt{file\_extent} (R2-A territory, chapter~13),
\texttt{snapshot\_metadata} and \texttt{snapshot\_name} (deferred), and
an "unsupported" bucket for record types the parser does not
recognise. A record falling into the unsupported bucket on a volume
that is otherwise in the allowlist triggers a typed validation gap
that the caller must consume.

\section{Topology Summary}

To enumerate one volume, the walker performs:

\begin{enumerate}
  \item Container OMAP lookup at scan XID for the volume's superblock
    OID, returning the volume superblock paddr.
  \item Volume superblock decode: the v1 feature allowlist is
    enforced; unknown incompatible bits, encryption, sealed-volume
    bits, and normalization-insensitive bits force
    \texttt{Unsupported(reason)}.
  \item Volume OMAP open; same encryption / unknown-flag gates as
    chapter~4.
  \item Volume OMAP lookup at scan XID for the FS-tree root virtual
    OID, returning the root paddr.
  \item FS-tree root validation as a virtual B-tree node.
  \item Recursive descent: at each internal node, an OMAP lookup per
    child entry; at each leaf node, decode of every record per the
    families above.
\end{enumerate}

Every reader in the project's Rust crate that touches an FS-tree node
goes through this path. There are no shortcuts; in particular there
is no path that treats an internal value as a paddr.

The next chapter handles what is inside the leaf records: the
variable-length tails, and the single source-backed rule for reading
them.
